\documentclass[../../../include/open-logic-section]{subfiles}

\begin{document}

\olfileid{sth}{replacement}{refproofs}
\olsection{附录：与替代公理相关的若干结果}

本节我们将在 $\ZF$ 中证明反射原理，并将证明在某种意义下反射原理等价于替代公理。此外，我们还将证明一个关于反射原理/替代公理之强度的有趣推论。\emph{警告：这可能是本教材中数学上最艰深的部分。}

我们先给出一个引理，为简洁起见，该引理采用\emph{上划线}记法来处理变元或对象的序列。即：“$\overline{a}_k$”是“$a_{k_1}$, \dots, $a_{k_n}$”的缩写，其中 $n$ 由上下文确定。

\begin{lem}\ollabel{lemreflection}
对每个 $1 \leq i \leq k$，令 $\phi_i(\overline{v}_i, x)$ 为一个公式。则对每个 $\alpha$，存在某个 $\beta > \alpha$，使得对任意 $\overline{a}_1,
\ldots, \overline{a}_k \in V_\beta$ 及每个 $1 \leq i \leq k$，有：
\[
	\exists x\phi_i(\overline{a}_i, x) \rightarrow (\exists x \in V_\beta) \phi_i(\overline{a}_i, x)
\]
\end{lem}

\begin{proof}
我们如下定义项 $\mu$：$\mu(\overline{a}_1, \ldots,
\overline{a}_k)$ 是满足以下所有条件句（对 $1 \leq i \leq k$）的最小阶段 $V$：
\[
\exists x\phi_i(\overline{a}_i, x) \rightarrow (\exists x \in V) \phi_i(\overline{a}_i, x))
\]
容易验证对所有 $\overline{a}_1, \ldots, \overline{a}_k$，$\mu(\overline{a}_1, \ldots, \overline{a}_k)$ 存在。现在，利用替代公理与递归定理，定义：
\begin{align*}
	S_0 & = V_{\alpha+1}\\
	S_{n+1} & = S_n \cup \bigcup
	\Setabs{\mu(\overline{a}_1, \ldots, \overline{a}_k)}
	{\overline{a}_1, \ldots, \overline{a}_k \in S_n} \\
	S &= \bigcup_{m < \omega} S_n.
\end{align*}
每个 $S_n$，因而 $S$ 本身，都是 $V_\alpha$ 之后的阶段。现取定 $\overline{a}_1$, \dots,~$\overline{a}_k \in S$；于是存在某个 $n < \omega$ 使得 $\overline{a}_1$, \dots, $\overline{a}_k \in S_n$。取定某个 $1 \leq i \leq k$，并假设 $\exists x
\phi_i(\overline{a}_i,x)$。由构造可知 $(\exists x \in \mu(\overline{a}_1,
\ldots, \overline{a}_k))\phi_i(\overline{a}_i, x)$，因此 $(\exists x \in S_{n+1})\phi_i(\overline{a}_i, x)$，进而 $(\exists x \in S)\phi_i(\overline{a}_i, x)$。所以 $S$ 就是我们的 $V_\beta$。
\end{proof}


\noindent
我们现在可以相当直接地证明 \olref[sth][replacement][ref]{reflectionschema}：

\begin{proof}[证明]
取定 $\alpha$。不失一般性，我们可假设 $\phi$ 的联结词仅为 $\exists$、$\lnot$ 和 $\land$（因为它们具有足够的表达力）。将 $\phi$ 的每个子公式按复杂度排列为 $\psi_1, \ldots, \psi_k$，使得 $\psi_k = \phi$。根据 \olref{lemreflection}，存在 $\beta > \alpha$，使得对任意 $\overline{a}_i \in V_\beta$ 和每个 $1 \leq i \leq k$：
\begin{align}\label{reflectionnicelybehaved}
	\exists x\psi_i(\overline{a}_i, x) \rightarrow
	(\exists x \in V_\beta) \psi_i(\overline{a}_i, x)\tag{*}
\end{align}
通过对 $\psi_i$ 的复杂度归纳，我们将证明对任意 $\overline{a}_i \in V_\beta$，有 $\psi_i(\overline{a}_i) \leftrightarrow
\psi_i^{V_\beta}(\overline{a}_i)$。
若 $\psi_i$ 是原子公式，这是平凡成立的。该双条件句也表明，当 $\psi_i$ 是满足此性质的子公式的否定或合取时，$\psi_i$ 自身也满足此性质。因此，唯一值得注意的情形涉及量词。取定 $\overline{a}_i \in V_\beta$；则有：
\begin{align*}
	(\exists x \psi_i(\overline{a}_i, x))^{V_\beta}
	&\text{ 当且仅当 }
	(\exists x \in V_\beta)\psi_i^{V_\beta}(\overline{a}_i, x)
	&&\text{据定义}\\
	&\text{ 当且仅当 }
	(\exists x \in V_\beta)\psi_i(\overline{a}_i,  x)
	&&\text{据归纳假设}\\
	&\text{ 当且仅当 }
	\exists x \psi_i(\overline{a}_i, x)
	&&\text{据 \eqref{reflectionnicelybehaved}}
\end{align*}
这就完成了归纳；由于 $\psi_k = \phi$，结果得证。
\end{proof}

我们已在 $\ZF$ 中证明了反射原理，证明本质上遵循了 \citet{Montague1961}。现在我们要在 $\Z$ 中证明反射原理蕴含替代公理。该证明遵循 \citet{Levy1960}，但做了简化。

由于我们是在 $\Z$ 中讨论，不能完全按上述形式呈现反射原理。毕竟，我们是用“$V_\alpha$”记法来表述反射原理的，而这在 $\Z$ 中无法定义（参见 \olref[sth][spine][zf]{sec}）。因此，我们改为给出一个看似更弱的反射原理表述，如下：

\begin{defish}
\emph{弱反射。} 对任意公式 $\phi$，存在一个传递集 $S$，使得 $0$、$1$ 以及 $\phi$ 的任意参数都是 $S$ 的元素，并且 $(\forall \overline{x} \in S)(\phi \liff
\phi^S)$。
\end{defish}

为了用它证明替代公理，我们首先遵循 \citet[first part of Theorem 2]{Levy1960}，证明我们可以同时“反射”两个公式：

\begin{lem}[在 $\Z + \text{弱反射}$ 中]\ollabel{lem:reflect}
对任意公式 $\psi, \chi$，存在一个传递集 $S$，使得 $0$ 和 $1$（以及公式的任意参数）是 $S$ 的元素，并且 $(\forall \overline{x} \in S)((\psi \liff \psi^S) \land (\chi
\liff \chi^S))$。
\end{lem}

\begin{proof}
令 $\phi$ 为公式 $(z = 0 \land \psi) \lor (z = 1 \land \chi)$。

此处我们用了一个缩写；我们应当将“$z = 0$”展开为“$\forall t\, t \notin z$”，将“$z =1$”展开为“$\forall s(s \in z
\liff \forall t\, t \notin s)$”。但由于 $0, 1 \in S$ 且 $S$ 是传递集，这些公式对 $S$ 是\emph{绝对的}；也就是说，无论是否将其量词限制于 $S$，它们都适用于同一个对象。\footnote{更形式化地说，令 $\xi$ 为这两个公式之一，则 $\xi(z) \liff \xi^S(z)$。}

由弱反射，存在某个适当的 $S$ 使得：
\begin{align*}
	(\forall z, \overline{x} \in S)(&\phi \liff \phi^S)\\
	\text{即 }(\forall z, \overline{x} \in S)(&((z = 0 \land \psi) \lor (z = 1 \land \chi)) \liff {}\\
	&\phantom{(}((z = 0 \land \psi) \lor (z = 1 \land \chi))^S)\\
	\text{即 }(\forall z, \overline{x} \in S)(&((z = 0 \land \psi) \lor (z = 1 \land \chi))\liff {}\\
	&\phantom{(}((z = 0 \land \psi^S) \lor (z = 1 \land \chi^S)))\\
	\text{即 }(\forall \overline{x} \in S)(&(\psi \liff \psi^S) \land (\chi \liff \chi^S))
\end{align*}
第二个断言蕴含第三个，是因为“$z = 0$”和“$z=1$”对 $S$ 是绝对的；第四个断言成立，是因为 $0 \neq 1$。
\end{proof}

\noindent 现在只需遵循并简化 \citet[Theorem 6]{Levy1960} 的证明，即可得到替代公理：

\begin{thm}[在 $\Z + \text{弱反射}$ 中]\label{thm:replacement}
对任意公式 $\phi(v,w)$ 和任意 $A$，若 $(\forall x \in A)\lexists![y][\phi(x,y)]$，则
$\Setabs{y}{(\exists x \in A)\phi(x,y}$ 存在。
\end{thm}

\begin{proof}
取定 $A$ 使得 $(\forall x \in A)\lexists![y][\phi(x,y)]$，并
定义公式：
\begin{align*}
	\psi &\text{ 是 } (\phi(x, z) \land A = A)\\
	\chi &\text{ 是 } \lexists[y][\phi(x, y)]
\end{align*}
利用 \olref{lem:reflect}，由于 $A$ 是 $\psi$ 的参数，故存在传递集~$S$，使得 $0, 1, A \in S$（连同任意其他参数），并且满足：
\[
	(\forall x,z \in S)((\psi \liff \psi^S) \land (\chi \liff \chi^S))
\]
因此特别地：
\begin{align*}
	(\forall  x, z \in S)(&\phi(x, z) \liff \phi^S(x, z))\\
	(\forall x \in S)(&\exists y\phi(x, y) \liff (\exists y \in S)\phi^S(x, y))
\end{align*}
结合这两点，并注意到由于 $A \in S$ 且 $S$ 是传递的，有 $A \subseteq S$：
\begin{align*}
	(\forall x \in A)(&\exists y\phi(x, y) \liff (\exists y \in S)\phi(x, y))
\end{align*}
现在 $(\forall x \in A)(\lexists![y \in S])\phi(x, y)$，因为
$(\forall x \in A)\lexists![y][\phi(x, y)]$。于是由分离公理即得
$\Setabs{y \in S}{(\exists x \in A) \phi(x, y)} = \Setabs{y}{(\exists
x \in A) \phi(x, y)}$。
\end{proof}

\end{document}